\bookchapter{第五十一章：界面開發——Jetpack Compose 入門}

\chapterintro{界面者，用戶所見之全也。Compose 未興之前，安卓界面以 XML 書之——繁瑣如填空之題；2021 年穀歌出 Jetpack Compose：逕以 Kotlin 代碼「畫」界面，所見即所得，改一行而預覽立新。至 2026 年，Compose 已為新項目之絕對主流。此章吾儕學「畫」之道。}

\booksection{一、Compose 之思：界面即函數}

Compose 核心之思，一言蔽之：界面 = 函數。凡可顯於屏者（文字、按鈕、圖片）皆一「可組合函數」（Composable），以 @Composable 誌之。嵌套之如積木，則界面成焉。

\begin{lstlisting}[language=Kotlin]
@Composable
fun Greeting(name: String) {
    Text(text = "汝好，$name！")     // 一"文字之積木"
}

// 嵌套而組：大積木套小積木
@Composable
fun HomeScreen() {
    Column {                          // 豎而列之
        Greeting(name = "張三")
        Greeting(name = "李四")
        Button(onClick = { /* 點之何為 */ }) {
            Text("點我")
        }
    }
}
\end{lstlisting}

慎視 @Composable 此「冠」——所以告編譯器曰：此乃界面積木也。積木可無限相套，此即 Compose 之全祕。

\booksection{二、基礎組件：文字、按鈕、輸入框、圖片}

Compose 之組件即「積木之塊」，至常見者四：

\begin{center}
\begin{tabularx}{\linewidth}{|X|X|X|}
\hline
\textbf{組件} & \textbf{所為何事} & \textbf{必備參數} \tabularnewline
\hline
Text & 顯文字 & text（所欲顯之字） \tabularnewline
Button & 按鈕 & onClick（點擊之回調） \tabularnewline
TextField & 輸入框 & value、onValueChange（輸入變化之回調） \tabularnewline
Image & 顯圖片 & painter（圖片之資源） \tabularnewline
\hline
\end{tabularx}
\end{center}

\begin{center}\small\color{gray}Compose 四大基礎組件\end{center}

\begin{lstlisting}[language=Kotlin]
@Composable
fun BasicWidgets() {
    Column {
        Text(text = "此乃一段文字")
        Button(onClick = { println("被點擊了") }) {
            Text("按鈕")
        }
        var input by remember { mutableStateOf("") }
        TextField(
            value = input,
            onValueChange = { input = it },
            label = { Text("輸入若干") }
        )
    }
}
\end{lstlisting}

輸入框彼兩行中現 remember \{ mutableStateOf(“”) \}——此即 Compose「狀態」之魔法，本章第五節詳言之。今姑識其面焉。

\booksection{三、布局三件：Column、Row、Box}

組件既有，何以擺置？記三「盒子」：

\begin{itemize}
  \item Column：縱而列之（如列隊而取食）；
  \item Row：橫而列之（如並肩而坐）；
  \item Box：疊而置之（如撲克之牌相壘，宜為角標、遮罩）。
\end{itemize}

\begin{lstlisting}[language=Kotlin]
Column {
    Text("第一行")
    Row {
        Text("左邊")
        Text("右邊")
    }
    Box {
        Text("底層")
        Text("疊於上之角標")
    }
}
\end{lstlisting}

欲布局愈美乎？與組件加 Modifier（修飾符）——猶為積木「上妝」也：

\begin{lstlisting}[language=Kotlin]
Text(
    text = "美之文字",
    modifier = Modifier
        .padding(16.dp)          // 四周留白
        .fillMaxWidth()          // 滿其寬
        .background(Color.Blue)  // 背景色
        .clickable { /* 點擊 */ }// 可點擊
)
// 鏈式調用：相繼而行，序有其義（先 padding 後背景，背景不含邊距）
\end{lstlisting}

\begin{tcbwarn}{慎之：Modifier 之序有考}
padding(16.dp).background(Blue) 與 background(Blue).padding(16.dp) 其效不同——前者「先留白後塗色」（背景不含邊距），後者「先塗色後留白」（背景含邊距）。新手最易於此畫出「怪異之白邊」。未定之時，易序而試，預覽立告汝以答。

\end{tcbwarn}

\booksection{四、狀態：使界面「活」}

靜態之界面無味也。欲鈕點而數增一、輸入框可打字，則須狀態（State）。Compose 之則：界面恆忠以映狀態，狀態一變，界面自重畫焉。

\begin{lstlisting}[language=Kotlin]
@Composable
fun Counter() {
    // remember + mutableStateOf：記此將變之值
    var count by remember { mutableStateOf(0) }

    Column {
        Text(text = "汝點了 $count 次")
        Button(onClick = { count++ }) {   // 改狀態
            Text("點我 +1")                // 界面自動刷新
        }
    }
}
\end{lstlisting}

三行代碼，得一活之計數器。此 Compose 勝舊式寫法之處：毋庸手動「刷新界面」，改數據足矣。記二關鍵函數：remember（使值於界面重組時不被清空）與 mutableStateOf（告 Compose 曰：此值變則當重畫）。

用戶點擊

狀態改變（count++）

Compose 感知到變化

自動重畫界面（重組）

\begin{center}\small\color{gray}圖 3：Compose 之「狀態 界面」單向數據流\end{center}

\booksection{五、主題與深色模式：一配而全局生效}

模板項目中之 MyFirstAppTheme \{ ... \} 即主題也。其本 Material 3 設計之範，助汝一統理顏色、字體、形狀。欲支深色模式，僅需三步：

\begin{enumerate}
  \item 於主題中定淺色/深色兩套之配色（模板已自生 ColorScheme）；
  \item 系統切至深色模式時，isSystemInDarkTheme() 自判之；
  \item 用主題所供之色（如 MaterialTheme.colorScheme.primary），勿硬寫色值。
\end{enumerate}

\begin{tcbinfo}{小識：Material 3 動態取色}
自 Android 12 起，「Material You」能使 App 自汝壁紙取色。Compose 模板默認支之（dynamicColor）。汝之 App 於不同用戶之機上將自易其裝，此現代安卓之浪漫也。

\end{tcbinfo}

\booksection{六、預覽、動畫與實時編輯：Compose 之「外掛」}

Compose 開發體驗之「外掛」有三，一用則不可復返：

\begin{itemize}
  \item @Preview 預覽：與 Composable 加 @Preview 之註，Android Studio 右側立渲染界面——改碼所見即所得，不必每番裝於機上；
  \item 實時編輯（Live Edit）：2026 年之 Android Studio 支於模擬器/真機上「熱更新」部分之改，若易文字顏色之屬小調，機上秒變；
  \item 動畫 API：animateFloatAsState、AnimatedVisibility 數行代碼即令諸元平滑而過——「高級之感」往往僅欠此數十行耳。
\end{itemize}

\begin{lstlisting}[language=Kotlin]
@Composable
fun AnimatedBox() {
    var big by remember { mutableStateOf(false) }
    // 值變則自平滑至標之
    val size by animateDpAsState(if (big) 120.dp else 60.dp)

    Box(Modifier
        .size(size)
        .background(MaterialTheme.colorScheme.primary)
        .clickable { big = !big }
    )
}
\end{lstlisting}

\begin{tcbinfo}{小識：先功能後動畫}
新手常犯之病，乃「界面尚未做對，遽糾結於動畫」。先使功能通焉，終乃一統加動效——此亦「簡單至上」之則：動畫者錦上添花也，功能者雪中送炭也。

\end{tcbinfo}

\begin{tcbpractice}{來動手：作一「心情按鈕」}
易模板項目為心情切換之器：一鈕，點之則於 喜 哀 怒 之間循環而切，文與背景之色隨之而變。提示：以 when 依狀態擇文案與色。既成以遺友玩之，可獲第一波「哇汝能作 App 矣」。

\end{tcbpractice}

\begin{tcbsummary}{本章小結}
\begin{itemize}
  \item Compose 之思：界面 = @Composable 函數，積木式相套
  \item 四大組件：Text / Button / TextField / Image
  \item 布局三件：Column 縱列、Row 橫列、Box 疊置
  \item Modifier 連環「上妝」：padding、fillMaxWidth、background、clickable——其序有意
  \item 狀態者魂也：remember + mutableStateOf，數據一變界面自重畫
  \item 主題用 Material 3 + 動態取色，深色模式一配而全局生效
\end{itemize}

\end{tcbsummary}
